\hmsubsection{Software Architecture}{FS}{}\label{sec:software-architecture}
The software running on the Raspberry Pi 5 is one Python application
split into five layers. Each layer is a Python module that does one
thing. The layers communicate through plain Python objects, since the
whole app runs in a single process.
The five layers follow the data flow of a pick-and-place cycle:
acquisition, perception, planning, communication, and supervision.
\subsubsection*{Acquisition}
The \texttt{OAKCamera} class wraps the DepthAI v3 API and gives the
rest of the code a \texttt{read()}/\texttt{release()} interface like
OpenCV's \texttt{VideoCapture}. This is the only place that knows
about the OAK-D S2. The rest of the pipeline can work with a recorded
video instead of the live camera, which is useful for offline tests.
The acquisition layer also gives access to the camera's factory
focal length, used by the distance estimation later in the pipeline.
The full behaviour is described in Section \ref{sec:acquisition}.
\subsubsection*{Perception}
The perception layer has two classes. \texttt{SimpleBallDetector} runs
the image processing pipeline once per frame. \texttt{BallTracker}
keeps a stable integer ID for each ball across frames. The output of
the layer is a list of tracked detections, each with a colour, a
centre in pixel coordinates, a radius, a confidence score, and an ID.
The full behaviour is described in
Sections \ref{sec:image-processing}--\ref{sec:tracking}.
\subsubsection*{Planning}
The \texttt{ArmIK} class converts a target position in workspace
centimetres into five Dynamixel goal positions. The class has no
state of its own; it is called once per cycle. The full behaviour is
described in Section \ref{sec:ik-solver}.
\subsubsection*{Communication}
The \texttt{SerialBridge} class owns the USB serial link to the
OpenRB-150. It exposes one method per command in the protocol, and
each method handles writes, reads, timeouts, and retries. The rest
of the host code can treat the arm as if it were an in-process object.
The full behaviour is described in Section \ref{sec:comm}.
\subsubsection*{Supervision}
The \texttt{CycleSupervisor} class is the top-level state machine.
It sequences each pick-and-place cycle through five states:
\texttt{IDLE}, \texttt{SCAN}, \texttt{PLAN}, \texttt{MOVE}, and \texttt{HOME}. Transitions are driven by events from the lower layers.
The supervisor is also the only place where calibration data is loaded
and where faults reported by the firmware are interpreted. Scan-pose
current limiting was evaluated during integration but is not part of
the deployed supervisory behaviour (Section \ref{sec:thermal}).
\subsubsection*{Why this split}
The split keeps each layer testable on its own. Perception can be run
against recorded frames without the camera. Planning is a pure
function, so we test it with round-trip and symmetry checks
(Section \ref{sec:fk-validation}). Communication can be tested with a
small Python script that sends every command and checks the reply
tokens. Supervision is the only layer that needs the full system
running.
The boundary between the host and the firmware is the USB serial
link. The boundary between the host and the camera is the DepthAI
API. Everything else is local to the Pi.
